\documentclass[12pt]{article}
\usepackage{amsmath}
\usepackage{amssymb}
\usepackage{hyperref}
\usepackage{graphicx}
\usepackage{booktabs}
\usepackage{geometry}
\geometry{margin=1in}

\title{DNS — Divergence-based Navigation System}
\author{Denis Schult\\ \texttt{schltdns@gmail.com}}
\date{April 2026}

\begin{document}

\maketitle

\begin{abstract}
Large Language Models (LLMs) often converge on similar outputs due to shared training data and alignment constraints. This convergence can mask systematic errors and hide epistemic uncertainty. We introduce the \textit{Divergence-based Navigation System (DNS)}, a domain-agnostic, human-centered protocol that treats structured disagreement between multiple LLMs as a primary epistemic signal. DNS maps, types, and quantifies divergence (Δdiv), enforces explicit falsification rules, and requires transparent human synthesis (P1–P6). A pilot benchmark across three domains (formal logic, applied systems, complex systems) demonstrates that semantic dispersion scales with epistemic complexity. DNS is not a truth machine — it is a divergence-sensitive evaluation layer for AI safety, model comparison, and decision-making under uncertainty. The method is open source (Apache 2.0), fully documented, and accompanied by case studies, a protocol specification, and a reproducible benchmark.
\end{abstract}

\section*{Motivation}
\begin{quote}
\textbf{“DNS emerged because curious people stop being curious when they fear social costs.”} – Claude
\end{quote}

\noindent DNS was developed to address a simple observation: in safety-critical and high-uncertainty domains, forced consensus hides risk. By systematically generating, mapping, and interpreting divergence across independent LLMs, DNS makes epistemic friction visible and actionable.

\section{What DNS Is}

DNS is a \textbf{divergence-sensitive epistemic protocol} that uses multi-model disagreement as a navigational signal for decision-making under uncertainty.

It enables:
\begin{itemize}
    \item structured analysis under uncertainty
    \item adversarial multi-model triangulation
    \item detection of blind spots, bias, and drift
    \item human-in-the-loop synthesis (P6)
    \item psychologically safe exploration (low-social-risk environment)
    \item quantifiable uncertainty signals (Δdiv)
\end{itemize}

DNS is the core method; GNS/TNS are domain-specific applications.

\section{Why Divergence > Consensus}

Consensus hides risk. Divergence reveals structure.

DNS illuminates:
\begin{itemize}
    \item what models cannot see about themselves (self-bias)
    \item what humans cannot see without them (cognitive blind spots)
    \item where systems disagree (epistemic fault lines)
\end{itemize}

Contradiction is treated as \textbf{primary data}, not as an error.

\section{Team Architecture v1.5}

DNS uses a fixed multi-agent role matrix:

\begin{center}
\begin{tabular}{llll}
\toprule
\textbf{Role} & \textbf{Agent} & \textbf{Responsibility} \\
\midrule
S1 – Kinetic-Agent & Grok & Real-time signals, fast sampling \\
S2 – Structural-Architect & ChatGPT & Documentation, logical framing \\
S3 – Synthesis-Bridge & Gemini & Strategic synthesis, long-term framing \\
S4a – Visualizer & Copilot & Data synthesis, operational tooling \\
S4b – Boundary-Marker & Claude & Underspecification detection, normative checks \\
Ω – Falsificator & DeepSeek & Stress-testing, falsification \\
S4 – Operator & Human & Final synthesis, justification, reflection \\
\bottomrule
\end{tabular}
\end{center}

\noindent Diagram: \texttt{dns\_architecture.png}

\section{Core Metrics: Divergence Delta (Δdiv)}

DNS measures \textbf{epistemic friction}, not correctness.

\[
\Delta_{div} = \frac{1}{n} \sum_{i=1}^{n} \frac{|M_i - \bar{M}|}{\sigma_{KRI}}
\]

Where:
\begin{itemize}
    \item \(M_i\): output vector of model \(i\)
    \item \(\bar{M}\): mean of all model outputs (consensus baseline)
    \item \(\sigma_{KRI}\): normalization factor (for geopolitical applications only; in generic DNS set to 1)
\end{itemize}

\subsection*{Reference Benchmark (pilot, n=small)}

\begin{center}
\begin{tabular}{lcc}
\toprule
\textbf{Domain} & \textbf{Δdiv (observed)} & \textbf{Epistemic Profile} \\
\midrule
Formal Logic & ~0.05 & Convergent \\
Applied Systems & ~0.40–0.55 & Structured Divergence \\
Complex Systems & ~0.65–0.80 & Contested / High Uncertainty \\
\bottomrule
\end{tabular}
\end{center}

Full benchmark: \texttt{BENCHMARK.md}.

\section{The DNS Protocol (P1–P6)}

\begin{enumerate}
    \item \textbf{P1 – Hypothesize:} Define 2–3 falsifiable hypotheses with measurable thresholds.
    \item \textbf{P2 – Thresholds:} Set explicit falsification criteria.
    \item \textbf{P3 – Triangulate:} Run identical prompts through all DNS agents.
    \item \textbf{P4 – Map Divergence:} Categorize contradictions into substantive, normative, didactic, and architectural types.
    \item \textbf{P5 – Operator Reflection:} Evaluate drift, bias, blind spots (three mandatory questions).
    \item \textbf{P6 – Human Synthesis:} The operator produces the final judgment and justifies every deviation.
\end{enumerate}

Full walkthrough: \texttt{HOW\_IT\_WORKS.md} and \texttt{EXAMPLE\_RUN.md}.

\section{Falsification Rules}

DNS is explicitly falsifiable. Four rules define when a hypothesis or synthesis is invalid:

\begin{enumerate}
    \item \textbf{Persistent Divergence:} After 3 rounds, divergence pattern unchanged → hypothesis inconclusive.
    \item \textbf{Refusal Rule:} Model refusal persists with substitute models → hypothesis underspecified (return to P1).
    \item \textbf{Human Justification Rule:} Operator cannot justify ignoring a model’s argument → synthesis invalid.
    \item \textbf{External Review Rule:} Human reviewer finds critical flaw without knowing model origins → return to P1/P3.
\end{enumerate}

\section{Limitations (Explicit)}

\begin{itemize}
    \item \textbf{No empirical validation:} The method is documented as a falsifiable hypothesis, not a proven causal model.
    \item \textbf{Operator dependence:} DNS is not an autopilot. Quality depends on the human operator’s reflection and domain knowledge.
    \item \textbf{Model bias:} All current LLMs share Western-aligned training data and alignment structures.
    \item \textbf{No quantification:} Δdiv is currently a qualitative heuristic.
    \item \textbf{Small sample size:} The benchmark is based on a pilot run (n=small).
\end{itemize}

\section{Case Studies}

DNS has been applied to three domains:
\begin{itemize}
    \item Cognitive Safety \& Neurodiversity
    \item Labor Market 2030
    \item European Gas Market Dynamics
\end{itemize}

\section{Related Work}

DNS complements existing frameworks:
\begin{itemize}
    \item Delphi Method
    \item AutoGen / CrewAI
    \item Semantic Triangulation
    \item Red Teaming
\end{itemize}

\section{License \& Availability}

\noindent This work is licensed under the Apache License 2.0.  
High-level system prompts and internal weighting matrices used during development are not included in this publication.

\noindent All documentation, benchmark, protocol, and case studies are available at:\\
\url{https://github.com/schltdns/divergence-navigation-system}

\section{Citation}

Schult, D. (2026). \textit{DNS — Divergence-based Navigation System: A divergence‑sensitive epistemic protocol for multi‑model evaluation under uncertainty.} GitHub: \url{https://github.com/schltdns/divergence-navigation-system}

\section*{Acknowledgments}

The author thanks the six AI systems used as investigator proxies (Grok, ChatGPT, Gemini, Copilot, Claude, DeepSeek) whose systematic disagreements provided the empirical material for this work.

\begin{thebibliography}{9}
\bibitem{denzin1970} Denzin, N. K. (1970). \textit{The Research Act}. Aldine.
\bibitem{popper1959} Popper, K. R. (1959). \textit{The Logic of Scientific Discovery}. Hutchinson.
\bibitem{dalkey1963} Dalkey, N., \& Helmer, O. (1963). The Delphi method. \textit{Management Science}.
\bibitem{wu2023} Wu, Q., et al. (2023). AutoGen. \textit{arXiv:2308.08155}.
\bibitem{du2023} Du, Y., et al. (2023). Multiagent debate. \textit{arXiv:2305.14325}.
\end{thebibliography}

\end{document}
